\chapter{Environmental and Sustainability Risk Assessment}
\label{chap:Environmental and Sustainability Risk Assessment}

%---------------------------- SECTION ----------------------------
\section{A Rapidly Changing Landscape}
\paragraph{Of all the areas of applied risk assessment covered in Part VI of this book, environmental and sustainability risk assessment is perhaps the one undergoing the most rapid and the most fundamental transformation. What was, a generation ago, primarily a concern for organisations with significant industrial operations — factories, mines, chemical plants, waste facilities — has become, in the space of a relatively few years, a mainstream compliance and governance obligation for organisations across virtually every sector.}
\paragraph{The drivers of this transformation are multiple and mutually reinforcing. Scientific consensus on the reality and urgency of climate change has hardened significantly. Investor and stakeholder expectations around environmental responsibility have shifted dramatically. And regulators across the world — responding to both the science and the shifting expectations — have introduced, and are continuing to introduce, a rapidly expanding body of requirements around environmental risk assessment, climate-related disclosure, and sustainability governance.}
\paragraph{For compliance and legal professionals, the practical implication is clear: environmental and sustainability risk assessment is no longer a niche concern or a matter of corporate social responsibility peripheral to the main compliance agenda. It is a mainstream regulatory obligation — one that is still developing, in some respects still uncertain in its precise requirements, but one that is moving in one direction only.}
\paragraph{This chapter explores the full landscape of environmental and sustainability risk assessment — from the established framework of environmental impact and operational risk assessment to the rapidly developing world of climate-related financial risk and ESG disclosure. It is designed to give compliance professionals a thorough grounding in what is required, what is developing, and how to think about building an approach that is both currently adequate and capable of evolving with the landscape.}

%---------------------------- SECTION ----------------------------
\section{The Environmental Risk Assessment Landscape — An Overview}
\paragraph{Environmental risk assessment encompasses a wide and diverse range of activities — from the assessment of risks to human health and ecosystems from industrial operations, to the assessment of climate-related financial risks to investment portfolios. Before examining specific areas in detail, it is useful to map the overall landscape.}
\paragraph{At the broadest level, environmental risk assessment can be understood as addressing two distinct but related categories of risk:}
\paragraph{\textbf{Risks that organisations pose to the environment} — the risks that an organisation's activities create for the natural environment, for ecosystems, for human health, and for communities. This is the traditional focus of environmental regulation — managing the pollution, the waste, the habitat destruction, and the resource consumption that industrial and commercial activity can generate.}
\paragraph{\textbf{Risks that environmental and climate change poses to organisations} — the risks that environmental and climate conditions create for organisations themselves — physical risks from extreme weather, flooding, and changing temperatures; transition risks from the shift to a lower-carbon economy; and liability risks from the legal and regulatory consequences of environmental harm.}
\paragraph{Both categories are increasingly subject to regulatory expectation, and both need to be addressed in a comprehensive environmental and sustainability risk assessment framework. We will explore each in turn.}

%---------------------------- SECTION ----------------------------
\section{Risks That Organisations Pose to the Environment}

\subsection{The Regulatory Foundation — Environmental Permitting and Impact Assessment}
\paragraph{The regulation of the environmental risks posed by organisations is well established in most developed jurisdictions — with a legal framework that predates climate-related regulation by several decades and that continues to impose direct, enforceable obligations on organisations across a wide range of sectors.}
\paragraph{\textbf{Environmental permitting} is the primary mechanism through which regulators control the environmental risks of significant industrial and commercial operations. In England and Wales, the \textbf{Environmental Permitting (England and Wales) Regulations 2016} provide the framework under which operators of activities with significant environmental impact — waste management facilities, industrial installations, water discharge activities, mining operations — are required to obtain and operate within the conditions of an environmental permit.}
\paragraph{The permit system is, at its core, a risk management framework. Permit conditions are designed around the environmental risks associated with the permitted activity — specifying the controls that must be in place, the monitoring that must be conducted, and the reporting that must be made to the regulator. Operators are required to maintain and demonstrate compliance with those conditions on an ongoing basis.}
\paragraph{For compliance professionals in organisations that operate under environmental permits — or that are considering activities that would require one — the permit framework creates direct and specific risk assessment obligations. The assessment of environmental risks is not merely a best-practice exercise but a precondition for legal operation.}
\paragraph{\textbf{Environmental Impact Assessment (EIA)} is a related but distinct requirement, applicable to certain categories of development and project activity — major infrastructure projects, significant industrial developments, large-scale land use changes — that are likely to have significant effects on the environment. An EIA requires the developer to identify, describe, and evaluate the likely significant environmental effects of the proposed project before it is approved — providing the consenting authority with the information needed to assess whether, and on what conditions, the development should proceed.}
\paragraph{For compliance professionals involved in major projects — whether in the organisation's own right or as advisers to clients — understanding the EIA obligation and ensuring that it is properly discharged is a significant element of the project governance framework.}

\subsection{Operational Environmental Risk Assessment}
\paragraph{Beyond the formal permit and EIA frameworks, organisations with significant environmental impacts are expected to maintain ongoing operational environmental risk assessments — systematic assessments of the environmental risks associated with their day-to-day operations.}
\paragraph{The key environmental risk categories for operational assessment include:}
\paragraph{\textbf{Air quality and emissions} — risks from the release of pollutants to atmosphere, including combustion products, process emissions, volatile organic compounds, and dust. Air quality regulation imposes specific limits on emissions from many industrial processes, and the risk assessment should consider both the regulatory compliance dimension and the potential harm to health and ecosystems from excess emissions.}
\paragraph{\textbf{Water quality and discharge} — risks from the release of pollutants to water — surface water, groundwater, and the sewerage system. Water pollution incidents are among the most significant environmental enforcement risks for many organisations, particularly those operating near watercourses or with processes that involve significant volumes of liquid materials.}
\paragraph{\textbf{Land contamination} — risks from the release of pollutants to land, including through spills, leaks, and improper disposal. Land contamination can create long-lasting and expensive remediation obligations, and organisations operating on historically industrial sites carry particular exposure in this area.}
\paragraph{\textbf{Waste management} — risks from the generation, handling, storage, transport, and disposal of waste. Waste regulation is complex and detailed, and non-compliance — particularly in relation to hazardous waste — carries significant enforcement risk.}
\paragraph{\textbf{Resource consumption} — whilst traditionally less prominent in regulatory frameworks than pollution control, the consumption of energy, water, and raw materials is increasingly a focus of environmental regulation and reporting obligation — particularly in the context of sustainability and climate-related requirements.}
\paragraph{\textbf{Biodiversity and habitat} — risks from the impact of operations on habitats, species, and ecosystems. Biodiversity obligations have historically been most relevant to development and land management activities, but are increasingly relevant to a broader range of organisations as biodiversity net gain requirements and nature-related reporting frameworks develop.}

\subsection{Environmental Liability — The Legal Consequences of Environmental Harm}
\paragraph{The legal consequences of environmental harm — whether from a pollution incident, a contamination legacy, or a failure to comply with permit conditions — can be significant and enduring.}
\paragraph{In the UK, the \textbf{Environmental Liability Directive} — implemented through the \textbf{Environmental Damage (Prevention and Remediation) (England) Regulations 2015} — imposes strict liability on operators for certain categories of environmental damage, including damage to protected species and habitats, water damage, and land contamination that creates a significant risk to human health. The strict liability standard means that the operator is liable for the cost of remediation regardless of fault — the fact that the damage was caused accidentally or despite reasonable precautions does not remove the liability.}
\paragraph{Beyond the Environmental Liability framework, civil liability for environmental harm — through common law nuisance, negligence, and Rylands v Fletcher — provides a basis for claims by those affected by environmental damage. The costs of defending and settling environmental liability claims can be substantial, and they are not always adequately covered by standard insurance policies without specific environmental liability cover.}
\paragraph{For compliance professionals, the environmental liability dimension reinforces the importance of proactive environmental risk assessment and management — the cost of prevention is almost always substantially less than the cost of remediation and litigation.}

%---------------------------- SECTION ----------------------------
\section{Climate-Related Financial Risk — The New Frontier}
\paragraph{Alongside the established framework of operational environmental risk assessment, the past decade has seen the emergence of a new and rapidly developing category of environmental risk assessment: \textbf{climate-related financial risk} — the assessment of the risks that climate change poses to organisations and to the financial system.}
\paragraph{This is the area where the environmental risk assessment landscape is changing most rapidly, where regulatory expectations are most actively developing, and where the implications for compliance professionals are most significant and most far-reaching.}

\subsection{The Two Categories of Climate-Related Financial Risk}
\paragraph{The analytical framework that has become most widely adopted for thinking about climate-related financial risk was developed by the \textbf{Task Force on Climate-related Financial Disclosures (TCFD)} — an industry-led initiative established by the Financial Stability Board in 2015 to develop consistent, comparable climate-related financial risk disclosure recommendations. The TCFD framework distinguishes between two fundamental categories of climate-related risk:}
\paragraph{\textbf{Physical risks} — the risks arising from the physical manifestations of climate change — changes in temperature, precipitation patterns, sea levels, and the frequency and severity of extreme weather events. Physical risks can be:}
\begin{itemize}
    \bitem{Acute}{arising from specific extreme weather events such as floods, storms, droughts, and wildfires}
    \bitem{Chronic}{arising from longer-term shifts in climate conditions such as rising average temperatures, changing precipitation patterns, and sea level rise.}
\end{itemize}
\paragraph{Physical risks can affect organisations directly — through damage to physical assets, disruption of operations, supply chain failures, and changes in resource availability — and indirectly — through their effects on the broader economy, on markets, and on the organisations and communities with which they interact.}
\paragraph{\textbf{Transition risks} — the risks arising from the process of adjusting to a lower-carbon economy. As society, markets, and regulators respond to climate change, organisations face risks from:}
\begin{itemize}
    \bitem{Policy and regulatory changes}{carbon pricing, emissions trading schemes, energy efficiency requirements, restrictions on high-carbon activities, and mandatory climate disclosure.}
    \bitem{Technology changes}{the disruption of existing business models and assets by low-carbon technologies, and the risk of being left behind by the pace of technological transition.}
    \bitem{Market changes}{shifts in customer and investor preferences away from high-carbon products, services, and business models.}
    \bitem{Reputational risks}{the risk of being perceived as insufficiently responsive to climate change, with consequential effects on stakeholder relationships, customer loyalty, and employee attraction and retention.}
\end{itemize}

\subsection{The TCFD Framework — Disclosure as a Risk Assessment Driver}
\paragraph{The TCFD recommendations — published in 2017 and since updated — have become the most widely adopted framework for climate-related financial risk disclosure globally. They recommend that organisations disclose information about their climate-related risks and opportunities across four thematic areas:}
\paragraph{\textbf{Governance} — the board's oversight of climate-related risks and opportunities, and management's role in assessing and managing them.}
\paragraph{\textbf{Strategy} — the actual and potential impacts of climate-related risks and opportunities on the organisation's businesses, strategy, and financial planning — including scenario analysis of the potential effects of different climate pathways.}
\paragraph{\textbf{Risk management} — the processes used to identify, assess, and manage climate-related risks, and how those processes are integrated into the organisation's overall risk management framework.}
\paragraph{\textbf{Metrics and targets} — the metrics used to assess climate-related risks and opportunities, and the targets used to manage them}
\paragraph{The TCFD framework is significant not just as a disclosure framework but as a risk assessment framework — because effective TCFD-aligned disclosure requires genuine, substantive risk assessment activity. An organisation cannot credibly disclose its climate-related risks if it has not genuinely identified and assessed them. The disclosure obligation therefore drives the risk assessment obligation, in the same way that DPIA requirements under the GDPR drive data protection risk assessment.}

\subsection{TCFD — From Voluntary to Mandatory}
\paragraph{Whilst the TCFD recommendations began as a voluntary framework, they have been progressively incorporated into mandatory regulatory requirements in an increasing number of jurisdictions.}
\paragraph{In the \textbf{United Kingdom}, TCFD-aligned disclosure has been made mandatory on a phased basis for a significant range of organisations — including UK-listed companies, large registered companies, limited liability partnerships, and major financial institutions. The UK's approach has been driven by the FCA, the PRA, and the Department for Business and Trade, and the requirements are continuing to develop.}
\paragraph{In the \textbf{European Union}, the \textbf{Corporate Sustainability Reporting Directive (CSRD)} — effective from 2024 and applying progressively to a widening range of companies — requires detailed sustainability reporting including climate-related risk and opportunity disclosure under the \textbf{European Sustainability Reporting Standards (ESRS)}. The CSRD represents a significant expansion and standardisation of sustainability reporting obligations across the EU and, through its third-country provisions, for non-EU companies with significant EU operations}
\paragraph{In the \textbf{United States}, the \textbf{Securities and Exchange Commission (SEC)} has introduced climate-related disclosure rules for public companies — requiring disclosure of material climate-related risks, greenhouse gas emissions, and the financial impacts of severe weather events and carbon offsets.}
\paragraph{For compliance professionals in organisations subject to any of these frameworks, the mandatory disclosure obligations create direct and specific risk assessment requirements — and the quality of the underlying risk assessment will be the primary determinant of the quality and credibility of the disclosure.}

\subsection{Scenario Analysis — The Core Analytical Tool for Climate Risk}
\paragraph{The TCFD framework specifically recommends the use of climate scenario analysis as a tool for understanding and disclosing climate-related risks. Scenario analysis — introduced in Chapter~\ref{chap:Risk Assessment Tools and Methodologies} as a general risk assessment tool — is particularly well suited to climate-related risk assessment for several reasons:}
\paragraph{Climate change is a long-term, uncertain, and path-dependent risk — the precise physical and transition risks an organisation faces depend significantly on which climate trajectory the world follows and how quickly. Point estimates of climate-related risk are therefore of limited value — what is needed is an understanding of how the risk profile changes under different plausible climate futures.}
\paragraph{Climate scenario analysis typically examines the organisation's risk profile under several contrasting scenarios — including a scenario consistent with limiting global warming to 1.5°C or 2°C (requiring rapid, deep decarbonisation and therefore high transition risk), and a scenario consistent with higher levels of warming (implying lower transition risk but higher physical risk). By comparing the organisation's risk profile across these scenarios, decision-makers can develop a clearer understanding of the full range of potential exposures and the robustness of their strategy and business model across different climate futures.}
\paragraph{In financial services, the \textbf{Bank of England's Climate Biennial Exploratory Scenario (CBES)} — conducted in 2021 — required major UK banks and insurers to assess their financial resilience under three climate scenarios. This exercise demonstrated both the practical feasibility of climate scenario analysis at scale and the significant analytical challenges it involves — particularly in relation to data quality, modelling capability, and the length of the time horizons required.}
\paragraph{For compliance professionals outside the financial services sector, climate scenario analysis is less well established in practice but is a clear direction of regulatory travel. Organisations that invest in developing their scenario analysis capability now are likely to be better placed when mandatory requirements arrive in their sector.}

%---------------------------- SECTION ----------------------------
\section{ESG Risk Assessment — Beyond Climate}
\paragraph{The focus on climate-related risk, whilst important, is only one dimension of the broader \textbf{Environmental, Social, and Governance (ESG)} risk landscape that is reshaping compliance obligations across every sector. A comprehensive sustainability risk assessment framework needs to address the full ESG spectrum.}

\subsection{Environmental — Beyond Climate}
\paragraph{The environmental dimension of ESG extends beyond climate change to encompass.}
\paragraph{\textbf{Biodiversity and nature} — the risks and impacts associated with the loss of biodiversity, the degradation of ecosystems, and the depletion of natural resources. The Taskforce on Nature-related Financial Disclosures (TNFD) — modelled on the TCFD — has developed a voluntary framework for nature-related risk disclosure that is expected to follow a similar trajectory from voluntary to mandatory as climate-related disclosure. For organisations with significant land use, water use, or supply chain dependencies on natural systems, nature-related risk is a material and developing compliance concern.}
\paragraph{\textbf{Water risk} — the risks associated with water scarcity, water quality, flooding, and the organisation's impact on freshwater systems. Water risk is particularly material for organisations in water-intensive sectors — food and beverage, agriculture, semiconductor manufacturing, power generation — and for those operating in water-stressed geographies.}
\paragraph{\textbf{Pollution and waste} — beyond the operational compliance dimension discussed earlier, the environmental dimension of ESG requires organisations to assess and manage their broader contribution to pollution and waste generation — including through their supply chains and the lifecycle impacts of their products.}
\paragraph{\textbf{Circular economy} — the risks associated with continued reliance on linear, take-make-dispose business models in an environment where regulatory and market expectations are shifting towards circularity — reducing, reusing, and recycling materials rather than disposing of them.}

\subsection{Social — The Human Dimension of Sustainability Risk}
\paragraph{The social dimension of ESG encompasses the risks and impacts associated with the organisation's relationships with people — employees, customers, communities, and workers in the supply chain. Key social risk areas include:}
\paragraph{\textbf{Human rights and modern slavery} — the risks that the organisation's activities or supply chains involve, or contribute to, human rights violations — including forced labour, child labour, unsafe working conditions, and other forms of modern slavery. The \textbf{Modern Slavery Act 2015} imposes specific reporting obligations on large organisations operating in the UK, and the broader human rights due diligence framework is developing rapidly at both national and international levels.}
\paragraph{\textbf{Supply chain labour standards} — the risks associated with poor labour conditions, inadequate health and safety, and unfair treatment of workers in the organisation's supply chain. These risks create both direct reputational and legal exposure and, increasingly, direct regulatory risk as human rights due diligence legislation develops.}
\paragraph{\textbf{Community impact} — the risks associated with the organisation's impact on the communities in which it operates — including economic impacts, cultural impacts, and the environmental and social effects of its operations on local populations.}
\paragraph{\textbf{Diversity, equity, and inclusion} — the risks associated with inadequate diversity and inclusion within the organisation — including regulatory and legal risks, talent and retention risks, and reputational risks from perceived failure to meet evolving societal expectations.}

\subsection{Governance — The Structural Foundation}
\paragraph{The governance dimension of ESG reflects the recognition that the quality of an organisation's governance structures — its board composition, its accountability mechanisms, its approach to executive remuneration, its transparency, and its ethical standards — is a primary determinant of how well it manages both ESG and conventional risks.}
\paragraph{For compliance professionals, governance is the most directly relevant dimension of ESG — because it is where risk management structures, accountability frameworks, and compliance oversight arrangements sit. An organisation with strong governance is better placed to manage all other risks effectively; one with weak governance is vulnerable regardless of the sophistication of its risk frameworks.}

%---------------------------- SECTION ----------------------------
\section{Building an Environmental and Sustainability Risk Assessment Framework}
\paragraph{With the full landscape of environmental and sustainability risk in view, what does a practical, adequate risk assessment framework in this area look like?}

\subsection{Integration with the Enterprise Risk Framework}
\paragraph{The first and most important structural principle is that environmental and sustainability risk assessment should be \textbf{integrated} with the organisation's overall enterprise risk framework — not treated as a separate, parallel process managed by a different team with different methodologies and different governance.}
\paragraph{Climate-related physical risks are financial risks and operational risks — they affect assets, supply chains, revenues, and costs. Transition risks are strategic, regulatory, and reputational risks. Social risks are conduct, reputational, and legal risks. Governance risks are oversight and accountability risks. All of these map directly to the risk taxonomy established in Chapter~\ref{chap:Types of Risk} — and they should be assessed, managed, and reported within the same framework as the organisation's other material risks.}
\paragraph{The separation of ESG risk into a standalone sustainability function — disconnected from the risk management, compliance, and finance functions — is a governance model that is rapidly becoming obsolete. Regulators increasingly expect to see ESG risk integrated into mainstream risk governance, not siloed in a corporate responsibility team.}

\subsection{Materiality Assessment — Knowing What Matters Most}
\paragraph{Given the breadth of the ESG risk landscape, \textbf{materiality assessment} — the process of identifying which ESG risks and opportunities are most significant for the specific organisation — is an essential starting point for any sustainability risk assessment framework.}
\paragraph{Materiality in the ESG context has traditionally been assessed through the concept of double materiality — a framework that considers both:}
\begin{itemize}
    \bitem{Financial materiality}{the degree to which ESG factors affect the organisation's financial performance and position.}
    \bitem{Impact materiality}{the degree to which the organisation's activities affect the environment and society.}
\end{itemize}
\paragraph{The CSRD and the ESRS adopt the double materiality framework explicitly — requiring organisations to assess and disclose both the financial impacts of sustainability risks on the business and the sustainability impacts of the business on the environment and society. This is a more demanding framework than pure financial materiality — it requires organisations to consider their broader impact, not just their own exposure.}
\paragraph{For organisations subject to CSRD, the materiality assessment is a formal process — with defined methodology, stakeholder engagement requirements, and documentation standards. For organisations not yet subject to mandatory sustainability reporting, a materiality assessment remains good practice — ensuring that the risk assessment framework focuses on the issues that genuinely matter rather than attempting to address every possible ESG topic with equal depth.}

\subsection{Scenario Analysis for Climate Risk}
\paragraph{As discussed above, climate scenario analysis is a central element of TCFD-aligned risk assessment and is increasingly embedded in mandatory disclosure frameworks. Building the capability to conduct credible climate scenario analysis — including the data infrastructure, the modelling capability, and the governance process — is a priority for organisations subject to or approaching TCFD-aligned disclosure requirements.}
\paragraph{The development of that capability takes time — and organisations that have not yet begun should not wait for full mandatory requirements before starting. A phased approach — beginning with qualitative scenario analysis using published scenarios and publicly available data, and progressively developing quantitative modelling capability — is both practical and credible.}

\subsection{Supply Chain and Value Chain Assessment}
\paragraph{For most organisations, the most significant environmental and social risks are not within their own direct operations but in their supply chains and value chains — the upstream suppliers and downstream customers and users whose activities collectively determine the organisation's total environmental and social impact.}
\paragraph{\textbf{Scope 3 greenhouse gas emissions} — the emissions associated with the supply chain and value chain rather than the organisation's own direct operations and energy use — typically represent the largest proportion of most organisations' total carbon footprint. Assessing and managing Scope 3 emissions is a significant analytical and governance challenge — requiring engagement with suppliers, customers, and partners who may be at very different stages of their own sustainability journey}
\paragraph{Similarly, assessing human rights and labour standards risks in supply chains requires visibility into supplier practices that many organisations currently lack. Building that visibility — through supplier questionnaires, audit programmes, and contractual requirements — is a developing area of compliance practice that the regulatory direction of travel is likely to accelerate.}

\subsection{Metrics and Targets — Making the Assessment Actionable}
\paragraph{Environmental and sustainability risk assessment is only complete when it is connected to \textbf{metrics and targets} — measurable indicators of performance and defined goals for improvement. Without these, risk assessment in this area remains descriptive rather than actionable.}
\paragraph{Key metrics commonly used in environmental and sustainability risk assessment and reporting include:}
\begin{itemize}
    \bitem{Greenhouse gas emissions}{Scope 1, 2, and 3, expressed in tonnes of CO2 equivalent.}
    \bitem{Energy consumption and intensity}{total energy use and energy use per unit of output.}
    \bitem{Water consumption and intensity}{particularly important for water-intensive sectors and water-stressed geographies.}
    \bitem{Waste generation and recycling rates}{including hazardous waste separately.}
    \bitem{Environmental incidents}{spills, permit breaches, and regulatory notices.}
    \bitem{Supply chain sustainability performance}{supplier assessment scores, proportion of supply chain covered by sustainability assessments.}
\end{itemize}
\paragraph{Targets — whether expressed as absolute reductions, intensity improvements, or alignment with science-based pathways — provide the goal against which the metrics are measured and the basis on which progress can be reported to stakeholders.}

%---------------------------- SECTION ----------------------------
\section{The Regulatory Direction of Travel}
\paragraph{Before closing this chapter, it is worth looking ahead — because the regulatory landscape in environmental and sustainability risk assessment is moving faster than in almost any other area covered in this book}
\paragraph{Several significant developments are either underway or clearly signalled:.}
\paragraph{\textbf{Mandatory nature-related disclosure} — following the trajectory of TCFD, the TNFD framework is likely to become mandatory for a significant range of organisations in the coming years. Nature-related risk assessment — biodiversity, water, land use will become a mainstream compliance obligation.}
\paragraph{\textbf{Human rights due diligence legislation} — the EU's \textbf{Corporate Sustainability Due Diligence Directive (CSDDD)} — which entered into force in 2024 — requires large companies operating in the EU to identify, prevent, mitigate, and account for human rights and environmental adverse impacts in their own operations and across their value chains. Similar legislation is developing in other jurisdictions. For compliance professionals in multinational organisations, the CSDDD represents a significant new compliance obligation with direct risk assessment implications.}
\paragraph{\textbf{Green taxonomy and sustainable finance} — the EU's \textbf{Taxonomy Regulation} classifies economic activities as environmentally sustainable according to defined technical criteria — creating a framework that affects both the labelling of financial products as sustainable and the assessment of environmental risks in investment and lending decisions.}
\paragraph{\textbf{Transition planning} — regulators in financial services and increasingly across other sectors are developing expectations around \textbf{transition planning} — the articulation of a credible, time-bound strategy for aligning the organisation's business model and activities with a net-zero economy. Transition planning is, at its core, a forward-looking risk and strategy assessment exercise, and it will become an increasingly prominent compliance obligation in the years ahead.}

%---------------------------- SECTION ----------------------------
\newpage
\section{Chapter Summary}
In this chapter we learned about the following key points:
\begin{table}[H]
  \centering
  \renewcommand{\arraystretch}{1.5}
  \begin{tabularx}{\textwidth}{ | >{\raggedright\arraybackslash}p{0.3\textwidth} 
								| >{\raggedright\arraybackslash}p{0.35\textwidth} 
                                | >{\raggedright\arraybackslash}X | }
    \hline
    \textbf{Risk Category} & \textbf{Key Regulatory Framework} & \textbf{Primary Obligation} \\
    \hline
    \textbf{Operational environmental risk} & Environmental Permitting Regulations; EIA Regulations & Environmental permits; impact assessments for significant developments\\
    \hline
    \textbf{Environmental liability} & Environmental Liability Directive; common law & Remediation of environmental damage; civil liability\\
    \hline
    \textbf{Climate — physical risk} & TCFD; CSRD; FCA/PRA expectations & Identify and disclose physical climate risks; scenario analysis\\
    \hline
    \textbf{Climate — transition risk} & TCFD; CSRD; Carbon budgets; emissions trading & Identify and disclose transition risks; transition planning\\
    \hline
    \textbf{Biodiversity and nature} & TNFD; Biodiversity Net Gain & Nature-related risk assessment and disclosure\\
    \hline
    \textbf{Human rights and supply chain} & Modern Slavery Act; CSDDD & Due diligence across value chain; modern slavery reporting\\
    \hline
    \textbf{Sustainability reporting} & CSRD; ESRS; UK sustainability disclosure standards & Double materiality assessment; integrated sustainability reporting\\
    \hline
  \end{tabularx}
  \caption{Environmental and Sustainability Risk Categories}
\end{table}

\begin{table}[H]
  \centering
  \renewcommand{\arraystretch}{1.5}
  \begin{tabularx}{\textwidth}{ | >{\raggedright\arraybackslash}p{0.35\textwidth} 
                                | >{\raggedright\arraybackslash}X | }
    \hline
    \textbf{Principle} & \textbf{Why It Matters} \\
    \hline
    \textbf{Integration with enterprise risk} & ESG risks are financial, strategic, and operational risks — they belong in the mainstream framework\\
    \hline
    \textbf{Double materiality} & Both financial impacts on the organisation and impacts on the environment and society must be assessed\\
    \hline
    \textbf{Scenario analysis for climate} & Point estimates are insufficient — scenario analysis is needed to understand the range of possible futures\\
    \hline
    \textbf{Value chain scope} & The most significant risks are often in the supply chain — Scope 3 and human rights due diligence require it\\
    \hline
    \textbf{Metrics and targets} & Assessment without measurement and goals is not actionable\\
    \hline
  \end{tabularx}
  \caption{Environmental and Sustainability Risk Assessment Key Principles}
\end{table}

\paragraph{Key takeaways:}
\begin{itemize}
  \item{Environmental and sustainability risk assessment is one of the most rapidly developing areas of compliance obligation — and the direction of travel is clear.}
  \item{The established framework of operational environmental risk assessment — permitting, EIA, liability — remains directly relevant and directly enforceable for organisations with significant environmental impacts.}
  \item{Climate-related financial risk assessment — encompassing both physical and transition risk — is becoming a mainstream compliance obligation across a widening range of sectors and organisation sizes.}
  \item{ESG risk extends significantly beyond climate — biodiversity, human rights, supply chain labour standards, and governance are all areas of developing regulatory expectation.}
  \item{Integration of ESG risk into the mainstream enterprise risk framework — rather than treating it as a separate sustainability agenda — is both better governance and increasingly a regulatory expectation.}
  \item{The regulatory direction of travel — mandatory TCFD, CSRD, CSDDD, and transition planning — makes early investment in environmental and sustainability risk assessment capability a sound governance and compliance decision.}
\end{itemize}

\barquote{In Chapter~\ref{chap:Financial and Commercial Risk Assessment}, we move from the environmental to the financial — exploring financial and commercial risk assessment, the quantitative and analytical frameworks used to assess financial risks, and the specific considerations that arise for compliance professionals working in or with financially regulated organisations.}